\documentclass[11pt,a4paper]{article}
\usepackage[margin=24mm,headheight=14pt]{geometry}
\usepackage{amsmath,amssymb,amsthm,mathtools,booktabs,array}
\usepackage{lmodern,microtype,fancyhdr}
\usepackage{fontspec,fancyvrb}
\usepackage[hidelinks]{hyperref}
\newfontfamily\leanappendixfont{DejaVuSansMono.ttf}
\newtheorem{theorem}{Theorem}
\newtheorem{lemma}[theorem]{Lemma}
\newtheorem{proposition}[theorem]{Proposition}
\newcommand{\R}{\mathbb R}
\newcommand{\nn}{\R_{\ge0}}
\newcommand{\tr}{\operatorname{tr}}
\newcommand{\diag}{\operatorname{diag}}
\newcommand{\gap}{\operatorname{gap}}
\newcommand{\rev}{\operatorname{rev}}
\newcommand{\code}[1]{\texttt{#1}}
\DeclareMathOperator{\im}{im}
\pagestyle{fancy}
\fancyhf{}
\fancyhead[L]{\small Collatz basin densities and canonical traces}
\fancyhead[R]{\small Lucas Valbuena}
\fancyfoot[C]{\thepage}
\renewcommand{\headrulewidth}{0.3pt}
\setlength{\parindent}{1.1em}
\setlength{\parskip}{3pt}
\setlength{\abovedisplayskip}{7pt plus 2pt minus 2pt}
\setlength{\belowdisplayskip}{7pt plus 2pt minus 2pt}
\hypersetup{pdftitle={Collatz basin densities and canonical traces from first-passage stabilization},pdfauthor={Lucas Valbuena},pdfsubject={Consequences of Tao's first-passage stabilization for logarithmic densities and critical traces}}

\begin{document}
\thispagestyle{plain}
\begin{center}
{\LARGE\bfseries Collatz basin densities and canonical traces\\[3pt]from first-passage stabilization\par}
\vspace{8pt}
{\large Lucas Valbuena\par}
\vspace{4pt}
{\normalsize 28 September 2026}
\end{center}
\vspace{2pt}
\begin{abstract}
We deduce from Tao's first-passage stabilization that every fixed Collatz
basin has a logarithmic density and that countable orbit labels have
limits in total variation. Two incommensurable scale parameters remove
the possible logarithmic oscillation left by one scale. Weighted basin
densities give a finite critical Green limit at every positive integer.
A first-hit time law then identifies this function with the Ces\`aro
and Abel limits of normalized stationary Syracuse cylinder masses.
For the largest cylinder mass we prove
\(2^{-k}\le\max_N\pi(N+3^k\mathbb Z_3)\le(2k+3)2^{-k}\),
which determines its optimal uniform exponential rate. The canonical
function is sublinear uniformly among the large ancestors of any fixed
root. Its global sublinearity on nonperiodic targets is equivalent to
logarithmic density one of the eventually periodic set. Additional results of Mazur upgrade the basin and label limits to natural
density and imply that the canonical function is positive exactly at
integers prime to three. With that input, global sublinearity is
equivalent to eventual periodicity of every orbit; this condition and
the exclusion of additional cycles remain unresolved. The accompanying
Lean development checks selected proof steps, but the main analytic
theorem chain has not yet been formalized.
\end{abstract}

\section{Introduction and notation}
\label{sec:canonical-traces}

We use Tao's first-passage stabilization theorem to study fixed integer
basins. The arguments in this paper are analytic proofs with selected steps
checked in Lean. We prove
existence of logarithmic basin densities, total-variation limits for
countable orbit labels, and equality of two canonical critical traces
at every positive integer. A final proposition characterizes the
global sublinearity condition that remains unresolved. We then use
separate results of Mazur to obtain natural-density limits and positivity
at every target prime to three.

Use the shortcut Collatz map
\[
T(x)=
\begin{cases}
x/2,&x\text{ even},\\
(3x+1)/2,&x\text{ odd}
\end{cases}
\qquad(x\in\mathbb Z_{>0}).
\]
Write
\[
 \mathcal B_N=\{q\ge1:T^j(q)=N\text{ for some }j\ge0\},
 \qquad
 H_E(X)=\sum_{\substack{q\le X\\q\in E}}\frac1q.
\]
The logarithmic density of \(E\), when it exists, is
\(d_{\log}(E)=\lim_{X\to\infty}H_E(X)/\log X\).
For odd sets we also use relative logarithmic density, whose
denominator is
\[
 H_{\mathrm{odd}}(X)=\frac12\log X+O(1).
\]
A target is \emph{periodic} if it lies on a finite cycle; a
nonperiodic target can still be eventually periodic. A
\emph{divergent orbit} never enters a cycle. Its integer states
are distinct and tend to infinity.
Two positive integers belong to the same \emph{tail component}
when their forward orbits eventually meet.

\subsection{Relation to the 2026 descent and formalization results}

The density, target, time, and observable must be distinguished when
comparing the present statements with recent work. Shaik's September
revision~\cite{Shaik2026}, version 4.0.0, proves natural-density descent
to
\[
 (\log n)^{A_{\rm FP}}U(n),\qquad
 A_{\rm FP}=\frac1{2(1-H_2(\log_3 2))},
\]
where \(H_2\) is binary entropy, for every prescribed positive
\(U(n)\to\infty\). The shortcut and raw
Collatz witnesses have a common landing, clocks centered at
\(2\log n/\lambda\) and \(3\log n/\lambda\), respectively, with
\(O(\sqrt{\log n\log\log n})\) errors, and a common height bound
\(nR(n)\) for every prescribed positive \(R(n)\to\infty\).
His fixed-exponent specialization with a balanced height ceiling attains
the joint exceptional exponent
\((1-H_2(\log_3 2))A-1/2\). Our fixed-target first-hit laws concern
membership in a specified basin and do not claim new drift constants.
The target in Shaik's theorem still tends to infinity; that theorem
does not itself supply a limit for the density of a specified basin.

Nashida's fixed-barrier preprint~\cite{NashidaBarrier2026} proves
\[
 \#\{n\le X:\operatorname{Col}_{\min}(n)>M\}
       \le KXM^{-c'},\qquad c'\simeq3.506\cdot10^{-6},
\]
uniformly in the endpoint \(X\) and integer barrier \(M\ge1\).
Its proof combines fixed-barrier counting assembled from Shaik's
formalization with Mazur's natural-density recursion. This is a
stronger quantitative tightness estimate than is needed below;
together with our existence of the minimum law it gives
\(\mu_{\min}((M,\infty))\le KM^{-c'}\).
It neither identifies a specified basin's mass nor makes a zero-density
exceptional set empty. Nashida's generalized-map
preprint~\cite{NashidaGeneral2026} proves the analogous power saving
for the Gon\c calves--Greenfeld--Madrid family. His separate
stopping-time preprint~\cite{NashidaTerras2026} improves the
Terras-type exponent for long stopping times. That count ranges over
starting integers and is distinct from the uniform packing bound for
one finite trajectory of distinct states used here.

Mazur's formalization of Tao~\cite{MazurTao2026} includes the
first-passage stabilization theorem, rather than only Tao's final
almost-boundedness conclusion. Its checked source fixes the scale
parameter at \(\alpha=1001/1000\). We rebuilt that theorem and
adapted its checked source to \(\alpha=2001/2000\), retaining the
real-threshold first-passage and rate conclusions. The two separate
Lean checks supply exactly the scales used in
Proposition~\ref{prop:ct-tao-input}; they do not assert a formally
checked continuum of scale parameters.
His natural-density transport theorem~\cite{MazurNatural2026}
compares uniform and logarithmic first-passage laws on the same
source block. Section~\ref{sec:ct-mazur-consequences} uses that
additional input to upgrade our basin and label limits to natural
density. His predecessor theorem~\cite{MazurPredecessors2026}
also implies positivity of the canonical trace at every integer
prime to three. These are credited consequences of those external
results, not claims that the external manuscripts state our
weighted or canonical-limit theorems.

The formal verification boundaries differ between these sources.
Shaik's version-4 theorem map distinguishes its checked principal
chain from additional manuscript-only refinements; the separate
Palomar version-3 record concerns the preceding revision.
Nashida's fixed-barrier theorem has no additional mathematical
Lean premise. His generalized natural-density theorem takes a
transcription of Matveev's bound as a hypothesis, and his
stopping-time improvement takes a transcription of Bugeaud's
bound; the cited published theorems supply those inputs to the
written proofs. The explicit numerical saving in the latter
paper is certified by rational interval calculations outside Lean.
Our inspection of those theorem statements and source records
does not constitute a local rebuild of their full formalizations.

\section{Uniform control of finite distinct trajectories}

For a finite shortcut path \(x_0,\ldots,x_\ell\), define its
correction product, excluding the terminal state, by
\[
 P(x_0,\ldots,x_\ell)
 =\prod_{\substack{0\le i<\ell\\x_i\ {\rm odd}}}
       \left(1+\frac1{3x_i}\right).
\]
If the path has \(r\) odd steps, then
\begin{equation}\label{eq:ct-product}
 x_\ell=\frac{3^r}{2^\ell}x_0P(x_0,\ldots,x_\ell).
\end{equation}
The bound below is uniform in the length and endpoints of a
distinct path.

\begin{lemma}\label{lem:ct-packing}
Put
\[
 \beta=H_2(\log_3 2)<1,\qquad
 H_2(t)=-t\log_2t-(1-t)\log_2(1-t).
\]
For every \(b\in(\beta,1)\), there are constants \(K_b,R_b<\infty\)
such that every finite path of distinct positive states satisfies
\begin{align}
 \#\{i:x_i\le X\}&\le K_bX^b\qquad(X\ge1),\label{eq:ct-packing}\\
 \sum_{\substack{i\\x_i\ge M}}\frac1{x_i}
     &\le R_bM^{b-1}\qquad(M\ge1).\label{eq:ct-packing-tail}
\end{align}
In particular there is a constant \(P_0<\infty\), independent of
the path, for which \(1\le P\le P_0\).
\end{lemma}
\begin{proof}
Choose \(\rho\in(1/2,\log_3 2)\) with \(H_2(\rho)=b\), and put
\(c=\rho\log_2 3<1\). Let
\(A(X)=\#\{i:x_i\le X\}\) and \(h=\lfloor\log_2X\rfloor\).
Length-\(h\) parity words correspond bijectively to residue classes
modulo \(2^h\), by successive determination of the parity at each
step. A word with \(r\) odd positions has affine slope
\(3^r/2^h\) and offset at most \((3/2)^r-1\). Hence a start at
most \(X\) with \(r\le\rho h\) has its \(h\)-step image below
\(3\cdot3^{\rho h}\le3X^c\).

Images whose complete \(h\)-step windows lie in the segment are
distinct. At most \(h\) indices are too close to its terminal
endpoint. The elementary binomial entropy bound gives at most
\(2^{hH_2(\rho)}\) words with more than \(\rho h\) odd positions,
and each residue class contains at most two possible starts in
\([1,X]\). Consequently
\[
 A(X)\le A(3X^c)+2X^b+\lfloor\log_2X\rfloor.
\]
This also holds if the segment has length less than \(h\).
Choose \(X_0\) so large that
\(3^bX^{-b(1-c)}\le1/2\) and \(3X^c<X\) for \(X\ge X_0\).
Since
\(\log_2X\le X^b/(be\log2)\), strong induction on integer \(X\),
using \(A(X)\le X\) in the initial range, proves
\eqref{eq:ct-packing} for a sufficiently large constant \(K_b\).
Taking integer parts handles real \(X\).

Divide the states at least \(M\) into the shells
\([2^jM,2^{j+1}M)\). Equation~\eqref{eq:ct-packing} gives
\[
 \sum_{x_i\ge M}\frac1{x_i}
 \le K_b2^bM^{b-1}\sum_{j\ge0}2^{j(b-1)}
 =:R_bM^{b-1}.
\]
Finally, \(\log P\le\frac13\sum_i1/x_i\), so one may take
\(P_0=\exp(R_b/3)\).
\end{proof}

The same tail estimate applies to any subset of the states.
In particular, if the odd source states of a distinct path all
exceed \(M\), its correction product is
\begin{equation}\label{eq:ct-prefix-product}
 1\le P\le\exp\!\left(\frac{R_b}{3}M^{b-1}\right).
\end{equation}
The intervening even states need not exceed \(M\).

\section{The first-passage input at two scales}

On positive odd integers let
\[
 S(q)=\frac{3q+1}{2^{v_2(3q+1)}},\qquad
 \lambda=\log(4/3),\qquad \delta=\frac{\lambda}{2}.
\]
Let \(\tau_x(q)\) be the first \(S\)-time at a value at most \(x\),
with value infinity on failure. Write \(\operatorname{Pass}_x(q)\)
for that landing value, and give it the default value \(1\) on
failure. For \(a>1\), let \(Q_y^{(a)}\) have logarithmic
distribution on the odd integers in \([y,y^a]\).
The default convention ensures the exact identity
\begin{equation}\label{eq:ct-pass-compose}
 \operatorname{Pass}_M\circ\operatorname{Pass}_x
 =\operatorname{Pass}_M\qquad(1\le M\le x),
\end{equation}
including trajectories on which passage fails.

\begin{proposition}[Tao's stabilization at the required scales]
\label{prop:ct-tao-input}
For each
\[
 a\in\left\{\frac{1001}{1000},\frac{2001}{2000}\right\}
\]
there are \(C,c>0\) such that, for all sufficiently large \(x\),
\begin{align}
 \Pr\bigl(\tau_x(Q_y^{(a)})=\infty\bigr)&\le Cx^{-c},
       &&y=x^a,\ x^{a^2},\label{eq:ct-tao-failure}\\
 d_{\rm TV}\!\left(
 \mathcal L(\operatorname{Pass}_x(Q_{x^a}^{(a)})),
 \mathcal L(\operatorname{Pass}_x(Q_{x^{a^2}}^{(a)}))
 \right)&\le C(\log x)^{-c}.\label{eq:ct-tao-tv}
\end{align}
\end{proposition}
\begin{proof}
For \(a=1.001\), these are
\cite[Proposition~1.11, equations~(1.19)--(1.20)]{Tao2022}.
We record the parameter check needed to use a second scale.
In the proof in \cite[Section~5]{Tao2022}, retain
\[
 n_0=\left\lfloor\frac{\log x}{10\log2}\right\rfloor,
 \qquad
 m_0=\left\lfloor\frac{a-1}{100}\log x\right\rfloor.
\]
The input Propositions~1.9 and~1.14 there are independent of \(a\).
All numerical margins in Section~5 remain strict, uniformly on the
compact interval \([1.0005,1.001]\) containing both scales.
Reduction of either source modulo
\(2^{3n_0}\) has error
\(O(x^{0.3-a}/\log x)\), smaller than the required \(x^{-0.3}\).
The short-time upper exponent used after equation~(5.5) is at most
\[
 a^3+\frac{\log_2 3-1.9}{10}<0.971501<0.99.
\]
The passage window in equation~(5.9) has endpoints of order
\((a-1)\log x/\lambda\) and
\((a^3-1)\log x/\lambda\). Its lower endpoint is much larger
than \(2m_0\), and its upper endpoint is below
\(0.01044\log x<n_0\). The buffers
\((\log x)^{0.6},(\log x)^{0.7},(\log x)^{0.8}\) are smaller
than these linear margins. The affine offset is at most
\(3^{n_0}=O(x^{(\log_2 3)/10})\), whereas the main term in the
regular-valuation regime is at least
\[
 x^{\,a+(\log_2(3/4))/10-o(1)}.
\]
Thus the relative offset estimates used in equations~(5.13)--(5.19)
also retain their margins.

For completeness, the residue summation in Lemma~5.3 can be bounded
uniformly without discarding the factor \(a_{m_0}\).
With \(A=a_1+\cdots+a_{m_0}\), the ordinary contribution is bounded
by a constant times \(2^{-A}a_{m_0}\), whose sum over all positive
tuples is \(2\). The additional large-total contribution is at most
\[
 3^n\Pr(G_1+\cdots+G_{m_0}>\tfrac12\log_2x),
 \qquad \Pr(G_i=j)=2^{-j}.
\]
Exponential Markov at parameter \((\log2)/2\) bounds it by
\[
 x^{(\log_2 3)/10-1/4}(1+\sqrt2)^{m_0}=O(x^{-\eta})
\]
for some \(\eta>0\), since \(m_0\le10^{-5}\log x\).
Finally \(m_0\asymp\log x\), the passage-window count and source
normalization have the same factor \((a-1)\log y\), and
\(n-m_0\ge m_0\). The application of Proposition~1.14 at depth
\(m_0\), and the remaining concentration errors in that proof,
therefore give the same negative power of \(\log x\).
This verifies both required scales by the published proof;
it is not an additional hypothesis about its printed proposition.
The accompanying replay of Mazur's Lean sources checks the two
scales separately. Its source law is the reciprocal-weighted law
on the inclusive odd interval \([y,y^a]\), and its first-passage
location uses the same default value \(1\) as above. Its distance
is full \(\ell^1\), so its bound implies
\eqref{eq:ct-tao-tv} with our total-variation convention.
\end{proof}

We use
\begin{equation}\label{eq:ct-two-scales}
 a=\frac{1001}{1000},\qquad b=\frac{2001}{2000}.
\end{equation}
Their logarithms have irrational ratio: equality of positive integer
powers would contradict the exponent of the prime \(7\).

\begin{lemma}[Two-scale averaging]\label{lem:ct-two-scale}
Let \(V\) be a Banach space and \(A:[t_0,\infty)\to V\) satisfy
\[
 \|A(t)\|\le Lt+C,\qquad
 \|A(t+h)-A(t)\|\le Lh+C\quad(h\ge0).
\]
For \(u>1\), put
\[
 p_u(t)=\frac{A(ut)-A(t)}{(u-1)t}.
\]
Suppose \(\log a/\log b\notin\mathbb Q\) and, uniformly as
\(t\to\infty\),
\[
 \|p_a(at)-p_a(t)\|=O(t^{-c})\quad(c>0),\qquad
 \|p_b(bt)-p_b(t)\|=o(1).
\]
Then \(A(t)/t\) converges in norm. Its error is
\(O(t^{-\eta})\) for every \(0<\eta<\min\{c,1\}\).
\end{lemma}
\begin{proof}
Set \(B(u)=e^{-u}A(e^u)\) and \(\ell=\log a\).
The increment bounds imply boundedness of \(B\) and
\[
 \|B(v)-B(u)\|\le2L(v-u)+O(e^{-u})\quad(v\ge u).
\]
For
\[
 Q(u)=\frac{aB(u+\ell)-B(u)}{a-1}
\]
the first scale estimate reads
\(\|Q(u+\ell)-Q(u)\|=O(e^{-cu})\).
The geometric tail is summable, so completeness gives
\[
 P(u)=\lim_{k\to\infty}Q(u+k\ell),\quad
 P(u+\ell)=P(u),\quad Q(u)-P(u)=O(e^{-cu}).
\]
Extend \(P\) periodically to all real \(u\). The exact recurrence
\[
 (B-P)(u+\ell)
 =a^{-1}(B-P)(u)+(1-a^{-1})(Q-P)(u)
\]
shows, by iteration from one fixed interval of length \(\ell\),
that \(B(u)-P(u)=O(e^{-\eta u})\) for every
\(\eta<\min\{c,1\}\). The modulus for \(B\), shifted by \(k\ell\)
and then passed to the limit, makes \(P\) Lipschitz.

Put \(m=\log b\). The second scale estimate gives
\[
 bB(u+2m)-(b+1)B(u+m)+B(u)\longrightarrow0.
\]
Replace \(u\) by \(u+k\ell\) and use the preceding convergence to
obtain the same expression with \(P\) equal to zero.
For \(D(u)=P(u+m)-P(u)\), this says \(bD(u+m)=D(u)\).
Backward iteration and boundedness force \(D=0\).
Thus \(P\) has the two incommensurable periods \(\ell,m\);
continuity makes it constant. The asserted convergence and rate follow.
\end{proof}

\section{Elementary local passage estimates}

The following estimates concern the full unstopped valuation prefix.
They do not assume independence after a hitting time.

\begin{lemma}\label{lem:ct-local}
Uniformly for \(a\in[1.0005,1.001]\), let \(Q\) have logarithmic
distribution on the odd integers in either
\([x^a,x^{a^2}]\) or \([x^{a^2},x^{a^3}]\).
There are \(C,c>0\) such that
\begin{align}
 \Pr\!\left(\tau_x(Q)=\infty\ \text{or}\
 \left|\tau_x(Q)-\frac{\log(Q/x)}{\lambda}\right|
       >C(\log x)^{3/5}\right)
 &\le C e^{-c(\log x)^{1/5}},\label{eq:ct-local-time}\\
 \Pr\bigl(\operatorname{Pass}_x(Q)\le\sqrt x\bigr)
 &\le Cx^{-c},\label{eq:ct-small-landing}\\
 \sup_{m\le x}\Pr\bigl(\operatorname{Pass}_x(Q)=m\bigr)
 &\le Cx^{-c}.\label{eq:ct-local-atoms}
\end{align}
The last two estimates include the default atom on passage failure.
\end{lemma}
\begin{proof}
Put
\[
 R=\lfloor0.1\log x\rfloor,\qquad
 B=\lfloor\tfrac12\log_2x\rfloor,\qquad
 E=(\log x)^{3/5}.
\]
For the actual trajectory let
\[
 a_i=v_2(3S^{i-1}(Q)+1),\qquad A_k=\sum_{i=1}^k a_i.
\]
Reduction of either source law modulo \(2^{B+1}\) differs in total
variation from the uniform odd-residue law by
\begin{equation}\label{eq:ct-residue-error}
 O\!\left(\frac{2^B}{x^a\log x}\right).
\end{equation}
Indeed the reciprocal sum in each residue class differs from its
integral by \(O(x^{-a})\), while the source harmonic mass is a
positive constant times \(\log x\), uniformly in the stated range.
A prescribed positive valuation tuple of total \(A_R\le B\)
occupies one odd residue class modulo \(2^{A_R+1}\); under the
uniform law its probability is \(2^{-A_R}\). These tuples are
disjoint. Their probabilities, and that of their complement, are
therefore the same as for independent geometric variables
\(\Pr(G_i=j)=2^{-j}\), up to \eqref{eq:ct-residue-error}.

The moment generating function
\[
 \mathbb E e^{tG_i}=\frac{e^t}{2-e^t}\qquad(t<\log2)
\]
gives
\[
 \Pr\left(\left|\sum_{i=1}^kG_i-2k\right|>E\right)
 \le2\exp\!\left(-c\min\{E^2/k,E\}\right).
\]
A union bound for \(k\le R\), and \(B/R\to5/\log2>7\), show
that, outside probability \(Ce^{-c(\log x)^{1/5}}\),
\begin{equation}\label{eq:ct-regular-prefix}
 A_R\le B,\qquad |A_k-2k|\le E\quad(0\le k\le R).
\end{equation}
Here the geometric upper tail \(A_R>B\) and the residue error are
polynomially small in \(x\).

The exact affine formula is
\[
 S^k(Q)=3^k2^{-A_k}Q+F_k,\qquad 0\le F_k\le3^k.
\]
On \eqref{eq:ct-regular-prefix}, its main term is at least
\(x^a(3/4)^R2^{-E}\), so the relative offset is at most
\[
 x^{0.1\log4-a}2^E\le x^{-4/5}
\]
for sufficiently large \(x\). Thus, uniformly for \(k\le R\),
\begin{equation}\label{eq:ct-local-log-path}
 \left|\log S^k(Q)-(\log Q-\lambda k)\right|\le E.
\end{equation}
The central crossing time satisfies
\[
 \frac{a-1}{\lambda}\log x
 \le \frac{\log(Q/x)}{\lambda}
 \le\frac{a^3-1}{\lambda}\log x
 <0.01044\log x.
\]
Bracketing it by an error \(2E/\lambda+1\) stays inside
\([0,R]\). Equation~\eqref{eq:ct-local-log-path} rules out an
earlier crossing and guarantees a crossing by the upper endpoint.
This proves \eqref{eq:ct-local-time}. It also bounds the landing
below by \(x\exp(-3E-\lambda)\).

For the polynomial estimates, we use a less precise good event.
If \(\tau_x(Q)>R\), the product identity, with all its first \(R\)
odd source states exceeding \(x\), implies
\[
 A_R\log2
 <(a^3-1)\log x+R\log3+\frac{R}{3x}
 \le1.8R\log2
\]
for large \(x\). The limiting upper bound for \(A_R/R\) is at most
\[
 \log_2 3+\frac{1.001^3-1}{0.1\log2}<1.63.
\]
Since \(B\ge6R\) eventually, the geometric lower and upper tails
and \eqref{eq:ct-residue-error} give
\begin{equation}\label{eq:ct-poly-good}
 \Pr(\tau_x(Q)>R)+\Pr(A_R>B)\le Cx^{-c}.
\end{equation}
On the complementary event, a landing at most \(\sqrt x\) would
require the last valuation, whose source is greater than \(x\),
to exceed \(\frac12\log_2x+\log_2 3>B\). This proves
\eqref{eq:ct-small-landing}, including failure.

Finally fix a landing value \(m\). On the event
\(\tau_x(Q)\le R,\ A_R\le B\), its stopped positive valuation
tuple has total at most \(B\). The affine formula shows that this
tuple and \(m\) determine at most one starting integer.
There are \(\sum_{A=1}^B2^{A-1}<2^B\) such tuples, allowing all
lengths. Each source atom has probability
\(O(x^{-a}/\log x)\). Adding \eqref{eq:ct-poly-good} proves
\eqref{eq:ct-local-atoms}.
\end{proof}

\section{Logarithmic laws for basins and countable labels}

\begin{proposition}\label{prop:ct-label-laws}
Let \(I\) be countable and let \(F\) assign a label in \(I\) to
each positive odd integer. Suppose that, for all sufficiently large
\(x\), successful first passage preserves its label:
\[
 F(q)=F(\operatorname{Pass}_x(q))
 \quad\text{whenever }\tau_x(q)<\infty.
\]
Then the law of \(F\) under logarithmically uniform odd starts in
\([1,X]\) converges in total variation to a probability on \(I\).
This applies to tail components, global orbit minima, and every
fixed first-passage location. Their limits also hold for all positive
starts, with the same probabilities, after applying \(F\) to the odd
part.
\end{proposition}
\begin{proof}
In \(V=\ell^1(I)\), define
\[
 A_F(t)=\sum_{\substack{q\le e^t\\q\ {\rm odd}}}
                       \frac{\delta_{F(q)}}q.
\]
Its norm is \(t/2+O(1)\), and its increments have norm
at most \(h/2+O(e^{-t})\). Push
\eqref{eq:ct-tao-tv} forward under \(F\); the discrepancy between
the original and first-passage labels is supported on passage
failure, controlled by \eqref{eq:ct-tao-failure}.
The harmonic mass of the odd interval \([e^t,e^{at}]\) is
\((a-1)t/2+O(e^{-t})\). With \(t=a\log x\), normalization gives
\[
 \|p_a(at)-p_a(t)\|_1=O(t^{-c})
\]
for each of the two scales in \eqref{eq:ct-two-scales}.
Lemma~\ref{lem:ct-two-scale} gives norm convergence of \(A_F(t)/t\).
The limiting vector is nonnegative and has norm \(1/2\); after
multiplication by two it is a probability. Norm convergence gives
total-variation convergence, uniformly over every subset of labels.
In particular there is no loss of mass among countably many labels.

For a component label, successful passage stays in the same
component. For the global orbit minimum, every odd state before
passage exceeds \(x\), while the landing is at most \(x\);
an even intermediate state is followed by an odd state no larger
than itself. The minimum is therefore preserved. For the label
\(F=\operatorname{Pass}_M\) at fixed \(M\), use
\eqref{eq:ct-pass-compose} for \(x\ge M\).

These labels, extended by taking the odd part, are invariant under
doubling. Write a positive start as \(2^hq\), \(q\) odd.
For any fixed cutoff \(H\), the corresponding harmonic sum of
label vectors is
\[
 \sum_{h=0}^H2^{-h}A_F(\log X-h\log2).
\]
The omitted starts are multiples of \(2^{H+1}\), with logarithmic
proportion tending to \(2^{-H-1}\). First pass \(X\to\infty\),
then \(H\to\infty\), to obtain the same limiting probability.
\end{proof}

Let \(\mu_{\rm tail}\) and \(\mu_{\min}\) denote the component
and minimum probabilities. In particular every union of tail
components has a logarithmic density, and
\begin{equation}\label{eq:ct-min-tight}
 d_{\log}\{q:\min_{j\ge0}T^j(q)>M\}
       =\mu_{\min}((M,\infty))\longrightarrow0.
\end{equation}
For fixed \(M\), let \(\Lambda_M\) denote the limiting law of
\(\operatorname{Pass}_M\) on odd starts. Equation
\eqref{eq:ct-pass-compose} implies
\[
 (\operatorname{Pass}_L)_*\Lambda_M=\Lambda_L\qquad(L\le M).
\]
These probabilities do not determine the mass of any specified
component or the probability of convergence to one.

\begin{proposition}\label{prop:ct-basin-density}
Every fixed shortcut basin \(\mathcal B_N\) has a logarithmic density.
\end{proposition}
\begin{proof}
On successful first passage with \(x\ge N\), its membership can
change only if the even target \(N\) is visited during the last
halving chain and then omitted by the Syracuse acceleration.
That event has landing \(\operatorname{oddpart}(N)\), and for
\(x\ge N^2\) it is included in
\(\operatorname{Pass}_x\le\sqrt x\).
Thus for both source laws in Proposition~\ref{prop:ct-tao-input},
the expectations of
\(\mathbf1_{\mathcal B_N}(Q)\) and
\(\mathbf1_{\mathcal B_N}(\operatorname{Pass}_x(Q))\)
differ by \(O(x^{-c})\), using failure and
\eqref{eq:ct-small-landing}. For an odd target there is no skipped
crossing. This argument also covers the default value for a
periodic target such as one.

Applying \eqref{eq:ct-tao-tv} to the bounded indicator gives the
two-scale estimates for
\[
 A_N(t)=\sum_{\substack{q\le e^t\\q\ {\rm odd},\,q\in\mathcal B_N}}
                         \frac1q.
\]
Lemma~\ref{lem:ct-two-scale} proves convergence of \(2A_N(t)/t\).
For a start \(2^hq\), membership agrees with that of its odd part
unless its initial halving chain hits \(N\). These exceptional
starts lie on the single ray
\(\{2^h\operatorname{oddpart}(N):h\ge0\}\), of finite reciprocal
sum. The finite-valuation decomposition in the preceding proof
therefore gives the all-positive density with the same value.
\end{proof}

\section{Weighted densities and the Green limit}

For \(q\in\mathcal B_N\), let \(P_N(q)\) be the correction product
on the path to its first visit to \(N\), and define
\[
 w_N(q)=P_N(q)^{-1}\quad(q\in\mathcal B_N),\qquad
 w_N(q)=0\quad(q\notin\mathcal B_N).
\]
The first-hit path is distinct, even when \(N\) is periodic.
Lemma~\ref{lem:ct-packing} gives
\begin{equation}\label{eq:ct-weight-bounds}
 P_0^{-1}\mathbf1_{\mathcal B_N}
          \le w_N\le\mathbf1_{\mathcal B_N}.
\end{equation}

\begin{proposition}\label{prop:ct-weighted-density}
For every \(N\ge1\), the limit
\[
 D_N=\lim_{X\to\infty}\frac1{\log X}
                    \sum_{q\le X}\frac{w_N(q)}q
\]
exists. The odd harmonic sum has limit \(D_N/2\), and
\begin{equation}\label{eq:ct-D-density}
 P_0^{-1}d_{\log}(\mathcal B_N)
       \le D_N\le d_{\log}(\mathcal B_N).
\end{equation}
\end{proposition}
\begin{proof}
Fix a packing exponent \(\gamma\in(\beta,1)\).
On successful passage below \(x\), outside a skipped even target,
first-hit paths concatenate:
\[
 w_N(q)=\frac{w_N(\operatorname{Pass}_x(q))}{P_{\rm prefix}}.
\]
This identity includes the case where both weights are zero.
The prefix is distinct and its odd source states exceed \(x\).
By \eqref{eq:ct-prefix-product},
\[
 |w_N(q)-w_N(\operatorname{Pass}_x(q))|
                         \le Cx^{\gamma-1}.
\]
Use the bound one on failure or a skipped crossing. For the two
source laws, their probabilities are \(O(x^{-c})\), as in
Proposition~\ref{prop:ct-basin-density}. Hence, for fixed \(N\)
and all sufficiently large \(x\),
\[
 |\mathbb E w_N(Q)-\mathbb E w_N(\operatorname{Pass}_x(Q))|
                    \le Cx^{\gamma-1}+Cx^{-c}.
\]
Combined with \eqref{eq:ct-tao-tv}, this proves the two-scale
estimates for
\[
 A(t)=\sum_{\substack{q\le e^t\\q\ {\rm odd}}}\frac{w_N(q)}q.
\]
This nonnegative scalar function satisfies the increment bounds in
Lemma~\ref{lem:ct-two-scale}, which therefore gives its limiting mean.

Initial halvings contribute no correction factor. Thus
\(w_N(2^hq)=w_N(q)\) outside the same single exceptional ray as
in the preceding proof. Its reciprocal sum is finite. The
finite-valuation decomposition gives the full mean \(D_N\) and
the factor \(1/2\) for the odd mean. Finally use
\eqref{eq:ct-weight-bounds}.
\end{proof}

For \(1<s<2\), define the positive inverse operator and its
Green sum by
\begin{align}
 (\mathcal R_s f)(v)
   &=2^{-s}f(2v)
      +(3/2)^s\mathbf1_{v\equiv2\ ({\rm mod}\ 3)}
                                f((2v-1)/3),\label{eq:ct-Rs}\\
 M_s(v)&=\sum_{j\ge0}\mathcal R_s^j\mathbf1(v),\qquad
 \kappa_s=2^{-s}+\frac{(3/2)^s}{3},\qquad
 C_s(v)=(1-\kappa_s)M_s(v).\label{eq:ct-Green-def}
\end{align}
Strict convexity gives \(\kappa_s<1\) on \(1<s<2\), and
\begin{equation}\label{eq:ct-kappa-expansion}
 1-\kappa_s=\delta(s-1)+O((s-1)^2).
\end{equation}
If \(N\) is periodic, let \(\ell_N\) be its shortcut period,
\(r_N\) its number of odd steps, and
\(\beta_N=3^{r_N}/2^{\ell_N}\). The return map has a strictly
positive affine offset, so \(0<\beta_N<1\). Put
\[
 B_N=
 \begin{cases}
 1,&N\text{ nonperiodic},\\
 (1-\beta_N)^{-1},&N\text{ periodic}.
 \end{cases}
\]
This scalar \(B_N\) is distinct from the basin \(\mathcal B_N\).

\begin{theorem}\label{thm:ct-Green}
At every positive integer the Green sum is finite for \(1<s<2\),
and its critical limit exists:
\begin{equation}\label{eq:ct-Green-value}
 g(N):=\lim_{s\downarrow1}C_s(N)=\delta B_N N D_N.
\end{equation}
It is a finite nonnegative function satisfying
\begin{equation}\label{eq:ct-critical-harmonic}
 g(v)=\frac12g(2v)
       +\frac32\mathbf1_{v\equiv2\ ({\rm mod}\ 3)}g((2v-1)/3).
\end{equation}
Moreover \(g(N)>0\) if and only if \(d_{\log}(\mathcal B_N)>0\).
\end{theorem}
\begin{proof}
The inverse word corresponding to the first hit from \(q\) has
weight \((N/(qP_N(q)))^s\), by \eqref{eq:ct-product}.
At a nonperiodic target no endpoint can hit it at two different
times. At a periodic target all later hits are obtained by complete
cycle traversals. Consequently
\begin{equation}\label{eq:ct-first-hit-Dirichlet}
 M_s(N)=
 \begin{cases}
 N^s\displaystyle\sum_{q\ge1}w_N(q)^s q^{-s},
                       &N\text{ nonperiodic},\\[4pt]
 \displaystyle\frac{N^s}{1-\beta_N^s}
          \sum_{q\ge1}w_N(q)^s q^{-s},
                       &N\text{ periodic}.
 \end{cases}
\end{equation}
This proves finiteness by comparison with \(\zeta(s)\).

Write \(A(t)=\sum_{q\le e^t}w_N(q)/q=D_Nt+o(t)\).
For \(\varepsilon=s-1>0\), Stieltjes integration by parts gives
\[
 \varepsilon\sum_{q\ge1}w_N(q)q^{-1-\varepsilon}
       =\varepsilon^2\int_0^\infty e^{-\varepsilon t}A(t)\,dt
       \longrightarrow D_N.
\]
The convention at zero includes a possible atom \(q=1\).
The bound \(A(t)\le t+O(1)\) justifies the change of variables
and dominated convergence. For \(0\le w\le1\),
\[
 0\le w-w^s
 =\int_1^s(-w^r\log w)\,dr\le\frac{s-1}{e}.
\]
Thus replacing \(w_N\) by \(w_N^s\) in the preceding Abelian limit
changes it by at most
\((s-1)^2\zeta(s)/e\to0\).
Equations~\eqref{eq:ct-kappa-expansion} and
\eqref{eq:ct-first-hit-Dirichlet} prove \eqref{eq:ct-Green-value}.

The resolvent equation is \(M_s=1+\mathcal R_sM_s\).
It has only two inverse branches at any fixed integer, so passage
to the finite pointwise limits gives
\eqref{eq:ct-critical-harmonic}. The positivity equivalence follows
from \eqref{eq:ct-D-density}.
\end{proof}

\section{Uniform bounds for backward basins}

\begin{proposition}\label{prop:ct-basin-boundary}
There is a function \(\varepsilon(M)=O((\log M)^{-c})\to0\)
such that
\begin{equation}\label{eq:ct-small-ancestor}
 \mathcal B_v\cap[1,M]=\varnothing
       \quad\Longrightarrow\quad
 d_{\log}(\mathcal B_v)\le\varepsilon(M).
\end{equation}
Consequently, for every fixed positive root \(n\),
\begin{equation}\label{eq:ct-fixed-basin-boundary}
 \lim_{\substack{v\to\infty\\v\in\mathcal B_n}}
 d_{\log}(\mathcal B_v)=0,\qquad
 \lim_{\substack{v\to\infty\\v\in\mathcal B_n}}\frac{g(v)}v=0.
\end{equation}
Both limits are uniform over the indicated targets.
\end{proposition}
\begin{proof}
Fix \(a=1.001\), suppress its superscript on \(Q_y\), and write
\[
 \nu_x=\mathcal L(\operatorname{Pass}_x(Q_{x^a})),\qquad
 x_j=M^{a^j}.
\]
Pushing \eqref{eq:ct-tao-tv} through \(\operatorname{Pass}_M\),
and using \eqref{eq:ct-pass-compose}, gives
\[
 d_{\rm TV}\bigl((\operatorname{Pass}_M)_*\nu_{x^{a}},
                 (\operatorname{Pass}_M)_*\nu_x\bigr)
             \le C(\log x)^{-c}\qquad(x\ge M).
\]
The sum along \(x_j\) is at most
\(C(\log M)^{-c}/(1-a^{-c})\). Combining it with the local
atom estimate \eqref{eq:ct-local-atoms} at \(M\) proves, uniformly
in \(j\) and \(m\le M\),
\[
 \Pr(\operatorname{Pass}_M(Q_{x_j^a})=m)
       \le C M^{-c}+C'(\log M)^{-c}=:\eta(M).
\]
For
\(E_{M,m}=\{q\text{ odd}:\operatorname{Pass}_M(q)=m\}\),
the source intervals \([x_j^a,x_j^{a^2}]\) tile
\([M^a,\infty)\), apart from shared endpoints.
A final partial interval ending at \(X\) is contained in its full
interval, whose right endpoint is at most \(X^a\).
Shared endpoints have finite total reciprocal mass. Hence
\[
 H_{E_{M,m}}(X)\le\eta(M)H_{\rm odd}(X^a)+O_M(1)
\]
with an error independent of \(m\), and the relative odd upper
logarithmic density is at most \(a\eta(M)\).

Now suppose \(\mathcal B_v\cap[1,M]=\varnothing\), let
\(u=\operatorname{oddpart}(v)\), and put
\(m_v=\operatorname{Pass}_M(u)\). The value \(m_v=1\) is allowed
when passage fails. Every odd start reaching \(v\) has no state
at most \(M\) before that hit. Its first odd landing at most \(M\)
is consequently the same as that of \(u\), or both passages fail.
For an arbitrary start whose initial halving chain already reaches
\(v\), its odd part is exactly \(u\). Therefore
\[
 \mathcal B_v\subseteq\{2^hq:h\ge0,\ q\in E_{M,m_v}\}.
\]
The dyadic closure of an odd set of relative upper logarithmic
density at most \(d\) has ordinary upper logarithmic density at most \(d\):
sum finitely many valuations and bound the omitted starts by
multiples of \(2^{H+1}\), then let \(H\to\infty\).
This proves \eqref{eq:ct-small-ancestor} with
\(\varepsilon(M)=a\eta(M)\).

For fixed \(q,n\), the intersection of the forward orbit of \(q\)
with \(\mathcal B_n\) is finite. If nonempty, that orbit reaches
\(n\); its prefix is finite, and a later return to \(\mathcal B_n\)
would make \(n\) periodic, leaving only one finite cycle.
Thus
\[
 F_{M,n}=\bigcup_{q\le M}
       \bigl(\{T^j(q):j\ge0\}\cap\mathcal B_n\bigr)
\]
is finite. Above its maximum, every \(v\in\mathcal B_n\) has
no ancestor at most \(M\), so
\eqref{eq:ct-small-ancestor} gives the first limit in
\eqref{eq:ct-fixed-basin-boundary}.
Only finitely many targets in \(\mathcal B_n\) can be periodic:
any such target and \(n\) lie on the same finite cycle.
For all the remaining targets,
\[
 \frac{g(v)}v=\delta D_v\le\delta d_{\log}(\mathcal B_v).
\]
This proves the second limit.
\end{proof}

\section{A first-hit law for the number of odd steps}

For \(q\in\mathcal B_N\), let \(k_N(q)\) count the odd source
states before the first shortcut visit to \(N\). For odd \(q,N\)
this is the first Syracuse hitting time. For an even target it
includes the last odd step even when the hit occurs before the end
of its halving chain.

\begin{proposition}\label{prop:ct-time-law}
For every fixed \(N\ge1\) and \(\epsilon>0\), the odd set
\begin{equation}\label{eq:ct-time-exception}
 \left\{q\in\mathcal B_N,\ q\ {\rm odd}:
   \left|k_N(q)-\frac{\log q}{\lambda}\right|
                      >\frac{\epsilon\log q}{\lambda}\right\}
\end{equation}
has logarithmic density zero.
\end{proposition}
\begin{proof}
Again fix \(a=1.001\). On the same reference source
\(Q_{x^a}\), apply \eqref{eq:ct-local-time} at thresholds
\(x\) and \(y=x^{1/a}\). Its interval is
\([x^a,x^{a^2}]=[y^{a^2},y^{a^3}]\), so both local estimates
apply. With probability at least
\(1-Ce^{-c(\log x)^{1/5}}\), both passages succeed and
\[
 \tau_y(Q_{x^a})-\tau_x(Q_{x^a})
       =\frac{\log x-\log y}{\lambda}
                       +O((\log x)^{3/5}).
\]
On this event the difference is exactly
\(\tau_y(\operatorname{Pass}_x(Q_{x^a}))\).
Thus the displayed stage estimate is an event about a single
landing variable with law \(\nu_x\).

Fix a large bottom threshold \(M\), put \(x_i=M^{a^i}\), and
start from \(Q=Q_{x_j^a}\). Telescoping
\eqref{eq:ct-tao-tv} and using \eqref{eq:ct-pass-compose} give
\begin{equation}\label{eq:ct-stage-TV}
 d_{\rm TV}\bigl(\mathcal L(\operatorname{Pass}_{x_i}(Q)),
                 \nu_{x_i}\bigr)
 \le C\sum_{r=i}^{j-1}(\log x_r)^{-c}
 \le C'(\log x_i)^{-c}\qquad(i\le j).
\end{equation}
At \(i=j\) the laws agree. A union bound for the stage events,
transferred under \eqref{eq:ct-stage-TV}, together with the local
top passage estimate, has total error at most
\[
 Ce^{-c(\log x_j)^{1/5}}
 +\sum_{i=1}^j
   \left(Ce^{-c(\log x_{i-1})^{1/5}}
                           +C'(\log x_i)^{-c}\right)
 \le C''(\log M)^{-c'}
\]
uniformly in \(j\). Only marginal total variation is needed,
since each stage event is a function of its own landing variable.
No independence between stages is asserted.

On the simultaneous good event, top success and success of every
stage permit exact time concatenation. The main terms telescope,
and
\[
 \sum_{i=0}^j(\log x_i)^{3/5}
       \le C_a(\log Q)^{3/5}.
\]
It follows that
\begin{equation}\label{eq:ct-global-time-M}
 \Pr\left(\tau_M(Q)=\infty\ \text{or}\
  \left|\tau_M(Q)-\frac{\log Q-\log M}{\lambda}\right|
                   >C_a(\log Q)^{3/5}\right)
       \le C''(\log M)^{-c'}.
\end{equation}
The default-valued identities were used only to compare landing
laws; actual time concatenation is performed on the successful event.

Take \(M\ge N\). Among the finitely many odd \(u\le M\) in
\(\mathcal B_N\), let \(R_{M,N}\) be the maximum of \(k_N(u)\),
or zero if there are none. For an odd \(q>M\) in \(\mathcal B_N\),
passage below \(M\) must succeed. If it occurs before the first
visit to \(N\), counts concatenate. If an even \(N\) is first
visited inside the final halving chain of that passage, its odd-step
count is already exactly \(\tau_M(q)\). Therefore
\begin{equation}\label{eq:ct-final-count}
 0\le k_N(q)-\tau_M(q)\le R_{M,N}.
\end{equation}
For fixed \(M,N,\epsilon\), the bounded terms
\(R_{M,N}+(\log M)/\lambda\) and the \(O((\log q)^{3/5})\) error
are smaller than \(\epsilon\log q/\lambda\) on all sufficiently
high source intervals. By \eqref{eq:ct-global-time-M},
the probability of \eqref{eq:ct-time-exception} on each such
interval is at most \(C''(\log M)^{-c'}\).

The intervals \([x_j^a,x_j^{a^2}]\) tile \([M^a,\infty)\).
As in Proposition~\ref{prop:ct-basin-boundary}, covering the last
partial interval costs at most the multiplicative factor \(a\)
in logarithmic mass. Thus the odd-relative upper density of
\eqref{eq:ct-time-exception} is at most
\(aC''(\log M)^{-c'}\). The set does not depend on \(M\);
letting \(M\to\infty\) proves the claim.
\end{proof}

The conclusion is an absolute zero-density statement, also when
\(\mathcal B_N\) has density zero. Multiplying its harmonic
summands by any weights in \([0,1]\) preserves it. No bound
uniform over all starting integers or all targets is asserted.

The corresponding lower-clock constraint is also consistent with
Inselmann's natural-density trajectory estimate
\cite[Theorem~1.10]{Inselmann2024}. The two-sided fixed-target law
above uses the transported upper-time estimate as well.

\begin{proposition}[Conversion of the first-hit clocks]
\label{prop:ct-clock-conversion}
Let \(\ell_N(q)\) be the shortcut first hitting time of \(N\).
For every fixed \(N\), in logarithmic density on all positive starts
in \(\mathcal B_N\),
\[
 k_N(q)=\frac{\log q}{\lambda}+o(\log q),\qquad
 \ell_N(q)=\frac{2\log q}{\lambda}+o(\log q).
\]
Precisely, for each positive error tolerance the set violating either
corresponding relative error bound has logarithmic density zero.
\end{proposition}
\begin{proof}
Write \(q=2^hu\) with \(u\) odd. Outside the single ray whose
initial halving chain hits \(N\), one has
\(k_N(q)=k_N(u)\) and \(\ell_N(q)=h+\ell_N(u)\).
The ray has finite reciprocal sum. For every fixed \(h\),
\(\log q-\log u=h\log2\) is bounded; hence
Proposition~\ref{prop:ct-time-law} gives the odd-clock conclusion
on that valuation class. Finitely many such classes suffice,
because the remaining starts have upper logarithmic density at most
\(2^{-H-1}\), which tends to zero with the cutoff.

The exact first-hit product formula is
\[
 \ell_N(q)\log2
 =k_N(q)\log3+\log q-\log N+\log P_N(q).
\]
Since \(1\le P_N(q)\le P_0\), its last two terms beyond
\(\log q\) are bounded for fixed \(N\).
The identity \(\log3+\lambda=\log4=2\log2\) gives the
shortcut-clock coefficient.
\end{proof}

\section{Stationary cylinders and their maximal mass}

We define the stationary measure explicitly so that its values at
ordinary integers require no choice of an almost-everywhere density.
Let \(G_1,G_2,\ldots\) be independent with
\(\Pr(G_i=a)=2^{-a}\), \(a\ge1\), and put
\begin{equation}\label{eq:ct-pi-series}
 X=\sum_{j\ge1}3^{j-1}2^{-(G_1+\cdots+G_j)}
       \quad\text{in }\mathbb Z_3,\qquad \pi=\mathcal L(X).
\end{equation}
The series converges \(3\)-adically. Its projection modulo \(3^k\)
is Tao's Syracuse random variable
\cite[Section~1.4, equations~(1.22)--(1.24)]{Tao2022},
after reversal of the independent valuation coordinates.
The memoryless geometric law gives
\[
 \pi=\tfrac12 H_*\pi+\tfrac12 J_*\pi,\qquad
 H(x)=x/2,\qquad J(x)=(3x+1)/2.
\]
For \(v\in\mathbb Z_3\), define
\begin{equation}\label{eq:ct-rho}
 \rho_k(v)=3^k\pi(v+3^k\mathbb Z_3),\qquad \rho_0(v)=1.
\end{equation}
The stationary equation gives, for \(k\ge1\),
\begin{equation}\label{eq:ct-cylinder-recursion}
 \rho_k(v)=\frac12\rho_k(2v)
       +\frac32\mathbf1_{v\equiv2\ ({\rm mod}\ 3)}
                           \rho_{k-1}((2v-1)/3).
\end{equation}
Indeed the inverse image under \(J\) is empty outside its unit
class, and in that class a depth-\(k\) cylinder pulls back to a
depth-\((k-1)\) cylinder.

Iterating the first term, its remainder tends to zero since
\(\rho_k\le3^k\). Thus for positive integer \(v\),
\begin{equation}\label{eq:ct-odd-blocks}
 \rho_k(v)=\mathcal S_1^k\mathbf1(v),\qquad
 (\mathcal S_s f)(v)
 =\sum_{\substack{a\ge1\\2^av\equiv1\ ({\rm mod}\ 3)}}
       (3/2^a)^s f((2^av-1)/3).
\end{equation}
Every integral argument in this formula is a positive odd integer.
An inverse word of depth \(k\), with exponents of total
\(A=a_1+\cdots+a_k\), has weight \(3^k2^{-A}\).
Read forward, it is a shortcut path with \(k\) odd steps and
total length \(A\). Distinct words with the same \(A\) have
distinct endpoints, because an endpoint determines its forward
path and hence its parity word.

\begin{proposition}\label{prop:ct-cylinder-max}
For \(k\ge1\), let
\[
 q_k=\max_{v\bmod3^k}\pi(v+3^k\mathbb Z_3).
\]
Then
\begin{equation}\label{eq:ct-cylinder-max}
 2^{-k}\le q_k\le(2k+3)2^{-k}.
\end{equation}
Consequently
\[
 \lim_{k\to\infty}\frac{-\log q_k}{k\log3}=\log_3 2.
\]
For every \(\theta<\log_3 2\) there is \(C_\theta\) such that
\(\pi(v+3^k\mathbb Z_3)\le C_\theta3^{-\theta k}\)
for all \(v,k\); no exponent larger than \(\log_3 2\) has
this property.
\end{proposition}
\begin{proof}
Choose a positive representative \(N\) of a cylinder. For fixed
\(k,A\), an admissible inverse word ends at an odd integer \(m\)
with
\[
 T^A(m)=N=\frac{3^k}{2^A}m+b,\qquad
 0\le b\le(3/2)^k-1.
\]
The offset bound follows by induction, since even steps divide
the current offset and odd steps replace it by \(3b/2+1/2\).
The possible endpoints therefore lie in an interval of length
strictly less than \(2^{A-k}\). Its odd integers number at most
\(2^{A-k-1}+1\). By the fixed-length injectivity just noted,
the contribution of this total \(A\) to the cylinder probability
\(3^{-k}\rho_k(N)\) is at most
\[
 2^{-A}(2^{A-k-1}+1).
\]
Summing \(k\le A<5k\) gives at most
\((2k+2)2^{-k}\). For the remaining words, discard admissibility.
Their total cylinder mass is at most
\[
 \Pr(G_1+\cdots+G_k\ge5k)
 \le(2/3)^{5k}\bigl(\mathbb E(3/2)^{G_1}\bigr)^k
 =(32/81)^k\le2^{-k}.
\]
This proves the upper bound.

On the event \(G_1=\cdots=G_k=1\), of probability \(2^{-k}\),
the projection of \eqref{eq:ct-pi-series} modulo \(3^k\) is
\[
 \sum_{j=1}^k3^{j-1}2^{-j}=(3/2)^k-1\equiv-1\pmod{3^k}.
\]
This proves the lower bound. Taking logarithms gives the limiting
exponent. The polynomial factor \(2k+3\) can be absorbed in
\(3^{(\log_3 2-\theta)k}\) for every strictly smaller \(\theta\).
The lower bound excludes any larger exponent. The bounds retain
a possible polynomial loss at the endpoint exponent.
\end{proof}

This estimate controls stationary cylinders uniformly. It does not
provide an integer-height estimate uniform in the depth \(k\).
To control depth averages at fixed integers, we use the endpoint
representation and the first-hit time law.

\section{Identification of the cylinder and Green limits}

For \(0<z<1\), set
\[
 \mathcal A_z(N)=(1-z)\sum_{k\ge0}z^k\rho_k(N).
\]
We first give the cutoff estimate that will control exceptional
first-hit times.

\begin{lemma}\label{lem:ct-endpoint-tail}
The total weight of inverse odd-block words of depths
\(1\le k\le K\) and endpoints greater than \(N2^{6K}\) is at most
\begin{equation}\label{eq:ct-endpoint-tail}
 \frac98(64/81)^K.
\end{equation}
Moreover \(\rho_k(N)\le5kN+(64/81)^k\) for \(k\ge1\), so
\(\mathcal A_z(N)\) is finite.
\end{lemma}
\begin{proof}
For a depth-\(k\) word of total length \(A\), its endpoint
satisfies \(m\le2^AN/3^k\). An endpoint greater than \(N2^{6K}\)
therefore requires \(A>6K\). Discarding admissibility and summing
over compositions gives total weight at most
\[
 \sum_{k=1}^K3^k\Pr(G_1+\cdots+G_k>6K)
 \le(2/3)^{6K}\sum_{k=1}^K9^k
 \le\frac98(64/81)^K.
\]
For a fixed depth \(k\), apply the same Markov estimate to
the complementary tail \(A\ge6k\); its bound is
\((64/81)^k\). For each
\(k\le A<6k\), endpoints are distinct positive integers below
\(2^AN/3^k\), all with weight \(3^k2^{-A}\). Their combined
weight is at most \(N\). There are \(5k\) such lengths, proving
the remaining assertions.
\end{proof}

\begin{theorem}\label{thm:ct-canonical-identification}
For every positive integer \(N\),
\begin{equation}\label{eq:ct-canonical-identification}
 \lim_{K\to\infty}\frac1K\sum_{k=0}^K\rho_k(N)
 =\lim_{z\uparrow1}\mathcal A_z(N)
 =\lim_{s\downarrow1}C_s(N)
 =g(N)=\delta B_NND_N.
\end{equation}
No convergence of the individual sequence \(\rho_k(N)\) is
asserted.
\end{theorem}
\begin{proof}
The inverse odd-block expansion has a precise interpretation in
terms of odd starting endpoints. A first-hit path from odd \(q\)
to \(N\) has depth \(k_N(q)\) and weight \(Nw_N(q)/q\).
For an even target the last block stops at \(N\) without necessarily
continuing to its odd part; its exponent still counts exactly the
shortcut steps in that block. Thus the representation remains exact.

Define the first-hit partial sum
\[
 F_N(K)=N\sum_{\substack{q\ {\rm odd},\,q\in\mathcal B_N\\
                       k_N(q)\le K}}\frac{w_N(q)}q.
\]
For a nonperiodic target its terms are distinct inverse endpoints.
For a periodic target it still selects only the first-hit word of
each endpoint. In both cases
Lemma~\ref{lem:ct-endpoint-tail} bounds its part with
\(q>N2^{6K}\) by \eqref{eq:ct-endpoint-tail}.
Within that cutoff the exceptional set in
Proposition~\ref{prop:ct-time-law} contributes \(o(K)\), since
its reciprocal sum is \(o(\log(N2^{6K}))\) and \(0\le w_N\le1\).

Outside that exceptional set, for \(0<\epsilon<1\),
\[
 (1-\epsilon)\frac{\log q}{\lambda}
       \le k_N(q)\le
 (1+\epsilon)\frac{\log q}{\lambda}.
\]
Hence all good basin endpoints below
\(\exp(\lambda K/(1+\epsilon))\) are counted, while every good
counted endpoint is below
\(\exp(\lambda K/(1-\epsilon))\).
By Proposition~\ref{prop:ct-weighted-density},
\[
 \sum_{\substack{q\le X\\q\ {\rm odd}}}\frac{w_N(q)}q
       =\frac{D_N}{2}\log X+o(\log X).
\]
The cutoff, the \(o(K)\) exceptional contribution, and this identity
give
\[
 \frac{\delta ND_N}{1+\epsilon}
 \le\liminf_{K\to\infty}\frac{F_N(K)}K
 \le\limsup_{K\to\infty}\frac{F_N(K)}K
 \le\frac{\delta ND_N}{1-\epsilon}.
\]
Letting \(\epsilon\downarrow0\) proves
\begin{equation}\label{eq:ct-first-hit-Cesaro}
 F_N(K)/K\longrightarrow\delta ND_N.
\end{equation}
Independently of the time law, the same cutoff and \(w_N\le1\)
give \(F_N(K)\le C_N(K+1)\).

Define \(e_N=\mathbf1_{2\mid N}\). If \(N\) is nonperiodic,
every odd endpoint hits it only once, and hence
\[
 \sum_{k=0}^K\rho_k(N)=e_N+F_N(K).
\]
For odd \(N\), the empty word is already counted by the endpoint
\(q=N\) in \(F_N\). For even \(N\), no odd endpoint has depth
zero, so \(e_N\) supplies the single term \(\rho_0(N)=1\).
If \(N\) is periodic, later visits repeat its cycle. Each odd
endpoint contributes at depths \(k_N(q)+r_Nj\), with weights
multiplied by \(\beta_N^j\). Thus
\[
 \sum_{k=0}^K\rho_k(N)
 =e_N+\sum_{j=0}^{\lfloor K/r_N\rfloor}
                   \beta_N^jF_N(K-r_Nj).
\]
The even depth-zero singleton is outside this sum: positive-depth
odd-block words always have odd endpoints.
The domination \(F_N(K)\le C_N(K+1)\) permits dominated
convergence against the geometric series after division by \(K\).
Equation~\eqref{eq:ct-first-hit-Cesaro} gives the Cesàro limit
\(\delta ND_N/(1-\beta_N)\).

In either case, summation by parts gives the elementary Abelian
implication from a convergent Cesàro mean to
\(\mathcal A_z(N)\to g(N)\). The value agrees with
Theorem~\ref{thm:ct-Green}, proving all of
\eqref{eq:ct-canonical-identification}.
\end{proof}

There is a useful check on the normalization. Every finite inverse
word is uniquely partitioned, in traversal order, as
\[
 D^{a_1-1}F\cdots D^{a_k-1}F D^j,\qquad
 a_i\ge1,\quad j\ge0,
\]
where \(D(v)=2v\) and \(F(v)=(2v-1)/3\) is used only when
integral. The last run contributes \(2^{-sj}\).
Summing these nonnegative finite-word weights by Tonelli, and then
the geometric series in \(j\), yields
\[
 M_s(N)=\frac1{1-2^{-s}}\sum_{k\ge0}\mathcal S_s^k\mathbf1(N).
\]
With \(\kappa'_s=3^{s-1}/(2^s-1)\), one has
\[
 C_s(N)=(1-\kappa'_s)\sum_{k\ge0}\mathcal S_s^k\mathbf1(N),
 \qquad 1-\kappa'_s=\lambda(s-1)+O((s-1)^2).
\]
This agrees with \eqref{eq:ct-canonical-identification}: odd
endpoints have harmonic mean \(D_N/2\), and
\(\lambda D_N/2=\delta D_N\). No limit at a moving integer
target was interchanged with a critical parameter limit.

\section{Forward component mass and the remaining global condition}

A tail component is called divergent if its orbits never enter a
cycle; otherwise it is the basin of a finite cycle. The next
proposition makes the remaining global estimate precise.

\begin{proposition}\label{prop:ct-forward-components}
Let \(\mathcal C\) be a divergent component and
\(n_j=T^j(n)\) any forward orbit in it. Then
\[
 d_{\log}(\mathcal B_{n_j})\longrightarrow\mu_{\rm tail}(\mathcal C).
\]
The normalized trace has a component-dependent limit \(L_{\mathcal C}\)
with
\begin{equation}\label{eq:ct-component-ratio}
 \frac{\delta}{P_0}\mu_{\rm tail}(\mathcal C)
       \le L_{\mathcal C}
       :=\lim_{j\to\infty}\frac{g(n_j)}{n_j}
       \le\delta\mu_{\rm tail}(\mathcal C).
\end{equation}
The following conditions are equivalent:
\begin{enumerate}
\item \(g(v)=o(v)\) as \(v\to\infty\) through nonperiodic targets;
\item every divergent component has \(\mu_{\rm tail}\)-mass zero;
\item the eventually-periodic set has logarithmic density one;
\item \(g\) vanishes on every divergent component.
\end{enumerate}
\end{proposition}
\begin{proof}
The basins \(\mathcal B_{n_j}\) increase and exhaust \(\mathcal C\),
but this alone would not justify continuity of logarithmic density.
Fix \(M\). Each of the finitely many \(u\le M\) in \(\mathcal C\)
meets the chosen spine. Thus, for some finite \(J(M)\), every such
\(u\) reaches \(n_j\) for \(j\ge J(M)\). It follows that
\[
 \mathcal C\setminus\mathcal B_{n_j}
       \subseteq\{q:\min_{i\ge0}T^i(q)>M\}
                    \qquad(j\ge J(M)).
\]
Equation~\eqref{eq:ct-min-tight} and the existence of the relevant
densities prove the asserted basin-density limit.

For \(q\in\mathcal C\), Lemma~\ref{lem:ct-packing} gives a
convergent whole-orbit product
\[
 P_\infty(q)=
 \prod_{\substack{i\ge0\\T^i(q)\ {\rm odd}}}
        \left(1+\frac1{3T^i(q)}\right),
 \qquad 1\le P_\infty(q)\le P_0.
\]
Set \(W_{\mathcal C}(q)=1/P_\infty(q)\) on \(\mathcal C\) and
zero elsewhere. On \(\mathcal B_{n_j}\), first-hit factorization
gives
\[
 W_{\mathcal C}(q)=\frac{w_{n_j}(q)}{P_\infty(n_j)},
             \qquad P_\infty(n_j)\longrightarrow1.
\]
The observables
\(W_{\mathcal C}\mathbf1_{\mathcal B_{n_j}}\) increase, and
their logarithmic means are \(D_{n_j}/P_\infty(n_j)\).
Their error from \(W_{\mathcal C}\) is bounded by
\(\mathbf1_{\mathcal C\setminus\mathcal B_{n_j}}\), whose
density tends to zero. Consequently the full weighted mean exists:
\[
 d_{\mathcal C}
 =\lim_{X\to\infty}\frac1{\log X}
       \sum_{\substack{q\le X\\q\in\mathcal C}}
                       \frac1{qP_\infty(q)}
 =\lim_{j\to\infty}D_{n_j}.
\]
Formula~\eqref{eq:ct-Green-value} gives
\(L_{\mathcal C}=\delta d_{\mathcal C}\), and the bounds on
\(P_\infty\) prove \eqref{eq:ct-component-ratio}.

Condition 1 implies 2 by applying
\eqref{eq:ct-component-ratio} along each divergent spine.
Conditions 2 and 3 are equivalent by countable additivity of
\(\mu_{\rm tail}\): bounded integer orbits are exactly the
eventually periodic ones, and they form the union of the periodic
components.

Assume condition 2 and fix \(\epsilon>0\). Choose finitely many
periodic components \(\mathcal C_1,\ldots,\mathcal C_r\) whose
total mass exceeds \(1-\epsilon\). For a cycle point \(c_i\) in
each one, \(\mathcal C_i=\mathcal B_{c_i}\).
The uniform fixed-basin limit
\eqref{eq:ct-fixed-basin-boundary} makes \(g(v)/v\) at most
\(\epsilon\) at all sufficiently large targets in their finite
union. At a nonperiodic target outside that union, its basin is
contained in the complement, so
\[
 \frac{g(v)}v\le\delta d_{\log}(\mathcal B_v)
 \le\delta\mu_{\rm tail}\!\left(
          \mathcal C_1^c\cap\cdots\cap\mathcal C_r^c\right)
 \le\delta\epsilon.
\]
This proves condition 1. Finally condition 2 forces density zero
for every basin in a divergent component, hence \(g=0\) there by
\eqref{eq:ct-Green-value}. Conversely a positive-mass divergent
component has positive limiting ratio in
\eqref{eq:ct-component-ratio}, so cannot satisfy condition 4.
\end{proof}

The restriction to nonperiodic targets avoids an unproved bound
on the cycle factors \((1-\beta_N)^{-1}\) across hypothetical
additional cycles. The equivalent global condition in
Proposition~\ref{prop:ct-forward-components} has not been proved.
The preceding Tao-based argument alone does not exclude nonempty
components of logarithmic density zero. The additional predecessor
input in the next section excludes that possibility and proves
positivity at every unit. It still does not prove the global condition
or rule out additional positive cycles.

\section{Consequences of Mazur's natural-counting and predecessor theorems}
\label{sec:ct-mazur-consequences}

This section uses two additional results of Mazur. His comparison of
natural and logarithmic first-passage laws upgrades the density limits
above to natural density. His positive lower bound for each eligible
predecessor set determines the support of the canonical trace. The
conclusions below are deductions from those inputs and the preceding
sections; the inputs are not attributed to Tao's published theorem.

Put \(\alpha=1001/1000\). For the odd block
\[
 B_y=\{q\ge1:q\text{ odd},\ y\le q\le y^\alpha\},
\]
let \(U_y\) be its uniform law and \(L_y\) its law proportional
to \(1/q\). We use the following consequence of
\cite[Theorem~3.1]{MazurNatural2026}. For every
\(0<d<5/143\), all sufficiently large \(x\), and either
\(y=x^\alpha\) or \(y=x^{\alpha^2}\), the block is nonempty and
\begin{align}
 \Pr\bigl(\tau_x(U_y)=\infty\bigr)
     &\le184x^{-1/32000},\label{eq:ct-mazur-natural-failure}\\
 \left\|\mathcal L(\operatorname{Pass}_x(U_y))
       -\mathcal L(\operatorname{Pass}_x(L_y))\right\|_1
     &\le384256(\log x)^{-d}.
       \label{eq:ct-mazur-natural-transport}
\end{align}
Here \(\|\cdot\|_1\) is full \(\ell^1\) distance, and first
passage has the same default value one as above. The cited theorem
actually bounds failure to hit by a specified finite time; the
never-hit bound \eqref{eq:ct-mazur-natural-failure} is weaker.

\begin{theorem}\label{thm:ct-natural-density-upgrade}
Every label covered by Proposition~\ref{prop:ct-label-laws} has the
same limiting law under uniformly chosen odd starts in \([1,X]\)
as under logarithmically chosen odd starts. Convergence is in total
variation. Extending the label by taking the odd part gives the same
limit under uniformly chosen positive starts.

Every fixed shortcut basin \(\mathcal B_N\) has natural density,
and
\[
 d_{\rm nat}(\mathcal B_N)=d_{\log}(\mathcal B_N),\qquad
 \lim_{X\to\infty}\frac1X\sum_{q\le X}w_N(q)=D_N.
\]
In particular, every union of tail components has natural density
equal to its logarithmic density. The natural laws of orbit minima
and of every fixed first-passage location also converge in total
variation to the probabilities already constructed above.
\end{theorem}
\begin{proof}
Let \(F\) be one of the labels in
Proposition~\ref{prop:ct-label-laws}, with values in a countable set
\(I\), and let \(\mu\) be its limiting probability. In
\(\ell^1(I)\), put
\[
 A_F(t)=\sum_{\substack{q\le e^t\\q\ {\rm odd}}}
                         \frac{\delta_{F(q)}}q.
\]
The proposition gives \(A_F(t)/t\to\mu/2\) in norm. Hence
\[
 \mathcal L(F(L_y))
 =\frac{A_F(\alpha\log y)-A_F(\log y)+O_{\ell^1}(y^{-1})}
        {\tfrac12(\alpha-1)\log y+O(y^{-1})}
 \longrightarrow\mu.
\]
The endpoint error accounts for a possible included lower endpoint.

On successful passage the label is unchanged. Coupling a source
with its own passage value therefore bounds the full \(\ell^1\)
distance between its source-label and passage-label laws by twice
its failure probability. Push
\eqref{eq:ct-mazur-natural-transport} forward under \(F\), and
use \eqref{eq:ct-mazur-natural-failure} and
Proposition~\ref{prop:ct-tao-input} for the two failure terms.
It follows that
\[
 \|\mathcal L(F(U_y))-\mathcal L(F(L_y))\|_1=o(1)
 \qquad(y=x^\alpha,\ x\to\infty).
\]
Thus \(\mathcal L(F(U_y))\to\mu\).

For an arbitrary large endpoint \(X\), take
\[
 x=X^{1/\alpha^2},\qquad y=X^{1/\alpha}.
\]
Then \(B_y\) is precisely the odd part of \([y,X]\).
The proportion of positive odds below \(y\) among those at most
\(X\) is \(O(X^{1/\alpha-1})=o(1)\). Replacing the uniform
law on this block by the uniform law on all positive odds at most
\(X\) changes full \(\ell^1\) distance by at most twice that
proportion. This proves the odd-label assertion.

For the basin and weighted assertions, fix \(N\). The only
additional obstruction to passage invariance of
\(\mathbf1_{\mathcal B_N}\) is a skipped crossing of an even
target. For \(x\ge N^2\), that event is contained in
\(\operatorname{Pass}_x\le\sqrt x\), as in the proof of
Proposition~\ref{prop:ct-basin-density}. Lemma~\ref{lem:ct-local}
gives probability \(O(x^{-c})\) under \(L_y\).
Transporting this event through
\eqref{eq:ct-mazur-natural-transport} gives probability \(o(1)\)
under \(U_y\) as well. Failure to pass also has probability
\(o(1)\) for both sources.

On successful passage without a skipped target, the first-hit
product factorization used in
Proposition~\ref{prop:ct-weighted-density} gives
\[
 w_N(q)=\frac{w_N(\operatorname{Pass}_x(q))}{P_{
                 q\to\operatorname{Pass}_x(q)}}.
\]
Every odd source before the passage is greater than \(x\), and
the successful prefix has distinct states. The tail estimate
\eqref{eq:ct-prefix-product} therefore implies
\[
 |w_N(q)-w_N(\operatorname{Pass}_x(q))|
       =O(x^{b-1})=o(1),
\]
where the fixed packing exponent satisfies \(b<1\).
Both observables lie in \([0,1]\). Applying
\eqref{eq:ct-mazur-natural-transport} to their passage values and
including the preceding exceptional probabilities proves
\[
 \mathbb E_{U_y}f-\mathbb E_{L_y}f=o(1),
 \qquad f=\mathbf1_{\mathcal B_N}\ \text{or}\ w_N.
\]
The logarithmic block means tend respectively to
\(d_{\log}(\mathcal B_N)\) and \(D_N\), by
Propositions~\ref{prop:ct-basin-density} and
\ref{prop:ct-weighted-density}, including their odd-relative
normalizations. The same top-block argument therefore proves these
limits for uniform odd starts.

It remains to pass to all positive starts. For labels extended by
the odd part, doubling leaves the label unchanged. For the two
scalar observables, doubling has the same property except when the
initial halving chain hits \(N\). All exceptional starts lie on
the ray \(\{2^jN:j\ge0\}\), which has only \(O_N(\log X)\)
members below \(X\). For each fixed \(h\), odd starts multiplied
by \(2^h\) have asymptotic proportion \(2^{-h-1}\) among all
positive starts, and their normalized laws or means have the limits
just proved. Sum over \(0\le h\le H\). The remaining starts
are multiples of \(2^{H+1}\), with proportion at most
\(2^{-H-1}+O(X^{-1})\). First let \(X\to\infty\), then
\(H\to\infty\). This proves the all-positive assertions and,
for labels, preserves convergence in full \(\ell^1\).
\end{proof}

The next conclusion uses Mazur's separate predecessor theorem
\cite[Theorem~1.1]{MazurPredecessors2026}. Write \(C\) for the
ordinary Collatz map, which sends an odd \(q\) to \(3q+1\),
and let
\[
 \mathcal P(a)=\{q\ge1:C^j(q)=a\text{ for some }j\ge0\}.
\]
The cited theorem states that, for every fixed \(a>0\) with
\(3\nmid a\), there are \(c_a>0\) and \(X_a\) such that
\[
 \#\{q<X:q\in\mathcal P(a)\}\ge c_aX
 \qquad(X\ge X_a,\ X\in\mathbb N).
\]
It assumes neither that the target reaches one nor that the target
is nonperiodic. Constants may depend on the target.

\begin{theorem}\label{thm:ct-positive-unit-support}
For every positive integer \(N\),
\[
 d_{\log}(\mathcal B_N)>0
 \quad\Longleftrightarrow\quad D_N>0
 \quad\Longleftrightarrow\quad g(N)>0
 \quad\Longleftrightarrow\quad 3\nmid N.
\]
The same equivalence holds with natural density in place of
logarithmic density. Every tail component has positive logarithmic
and natural density. Consequently the equivalent global conditions
in Proposition~\ref{prop:ct-forward-components} are also equivalent
to every positive orbit being eventually periodic.
\end{theorem}
\begin{proof}
The map comparison
\cite[equation~(1.3)]{MazurPredecessors2026} is
\[
 \mathcal P(2N)\subseteq\mathcal B_N\subseteq\mathcal P(N).
\]
For completeness, a shortcut orbit retains the ordinary orbit
except for the intermediate even output immediately after an odd
step. If an ordinary occurrence of \(2N\) is retained, its next
shortcut step reaches \(N\). If it is omitted, that combined
shortcut step itself ends at \(N\). This proves the first
inclusion; expanding shortcut steps proves the second.

If \(3\nmid N\), apply Mazur's theorem to \(2N\). The first
inclusion gives positive lower natural density of \(\mathcal B_N\).
It also gives positive lower logarithmic density: writing
\(A(t)=\#\{q\le t:q\in\mathcal B_N\}\), the count bound
gives \(A(t)\ge c_{2N}t-O(1)\) for large real \(t\), and
summation by parts yields
\[
 \sum_{\substack{q\le X\\q\in\mathcal B_N}}\frac1q
 =\frac{A(X)}X+\int_1^X\frac{A(t)}{t^2}\,dt
 \ge c_{2N}\log X-O_N(1).
\]
Thus the existing logarithmic density is positive, independently
of the natural-density upgrade. The bounds
\eqref{eq:ct-weight-bounds} and the Green formula
\eqref{eq:ct-Green-value} give \(D_N>0\) and \(g(N)>0\).

If \(3\mid N\), the only shortcut predecessor of any positive
multiple of three is its double: an odd shortcut step has output
congruent to two modulo three. Inducting backward along a finite
path gives
\[
 \mathcal B_N=\{2^jN:j\ge0\}.
\]
Its natural and logarithmic densities are zero. Hence \(D_N=0\)
and \(g(N)=0\), again by the weighted bounds and Green formula.
Theorem~\ref{thm:ct-natural-density-upgrade} supplies the asserted
natural-density equivalence.

Every positive orbit reaches an integer prime to three. Indeed,
remove its initial powers of two; an odd multiple of three has
its next shortcut value prime to three. A tail component containing
such a value \(N\) also contains \(\mathcal B_N\), so its
logarithmic density is positive. Its natural density is the same
by Theorem~\ref{thm:ct-natural-density-upgrade}.

One of the equivalent conditions in
Proposition~\ref{prop:ct-forward-components} is that every
divergent component have zero logarithmic density. Since every
component has positive density, this now means that there are no
divergent components. That is exactly eventual periodicity of all
positive orbits. Conversely, if every orbit is eventually periodic,
the eventually periodic set is all positive integers, so the
conditions of that proposition hold.
\end{proof}

The positivity of the trace on units thus follows after adding
Mazur's predecessor theorem. Global sublinearity remains unproved;
with this additional input it is equivalent to excluding divergent
orbits. Even that conclusion would still permit cycles other than
the known cycle through one. Excluding such cycles, and hence
proving the Collatz conjecture, is not achieved here.

\section{Formal verification scope}

The accompanying Lean development verifies Lemma~\ref{lem:ct-two-scale}
in its original Banach-valued form, including the convergence rate,
and checks the irrational ratio of the two numerical scales.
The trajectory development proves the parity count, entropy bound,
and actual counting recurrence, thereby closing the uniform packing
and reciprocal-tail estimates in Lemma~\ref{lem:ct-packing}.
It also proves the uniform finite correction-product bound, infinite
product convergence on distinct orbits, and the forward-tail limit.
The cylinder development defines the full countable arithmetic word
law, proves its expansion and identification with the explicit finite
Syracuse residue law, and checks the upper, lower, and maximum mass
bounds. Identifying those finite laws with the infinite stationary
probability measure on \(\mathbb Z_3\) remains a separate obligation.
The raw-to-shortcut
map inclusion used in Theorem~\ref{thm:ct-positive-unit-support} and
the predecessor description for multiples of three are also checked.

These checks do not yet formalize the complete analytic theorem chain.
Separate Lean~4.30.0-rc2 replays check Mazur's real-threshold
first-passage rate theorem at each required scale, and his original
natural-density transport root. The local Lean~4.27.0 development
does not yet connect these inputs to the density-existence arguments,
Green and cylinder limit identification, or the additional Mazur
interfaces. The repository's
\texttt{paper/basin-lean-claims.md} records the exact declarations and
remaining obligations. Its verifier builds the local proof files afresh
and audits their axiom dependencies separately from the matrix and
synchronization papers. The present batch comprises thirty-five modules
and 251 public theorem and lemma declarations; the axiom audit also
covers fifty-five definitions. The external replay records state
their separate source and dependency pins.

\section*{Acknowledgments}
The main first-passage input and stationary Syracuse law are due to
Terence Tao. The natural-density and positivity deductions use the
separate results of Lech Mazur. The related-work discussion credits
Idris Ali Shaik, Hiroyuki Nashida, and Manuel Inselmann for their
descent and stopping-time results.
The formal parity-vector counting module adapts M.~Sharpe's Lean
proofs under their MIT license; its source retains the attribution
and license notice.

\section*{Use of AI}
OpenAI Codex was used extensively for mathematical exploration and
preparation and revision of this manuscript. The proofs use the
explicitly cited analytic inputs. Selected proof steps have been checked
in Lean; the full analytic theorem chain has not yet been formalized.
Automated review does not constitute external peer review or endorsement
by the authors of the cited works.

\section*{License}
Copyright \textcopyright\ 2026 Lucas Valbuena.
This paper and its LaTeX source are licensed under
\href{https://creativecommons.org/licenses/by/4.0/}{Creative Commons
Attribution 4.0 International (CC BY 4.0)}. The original accompanying
code is licensed under MIT. Imported proof sources retain their own
licenses; the Mazur replay in the repository uses Apache~2.0.

\begin{thebibliography}{99}\small\raggedright
\bibitem{Tao2022} T. Tao.
\emph{Almost all orbits of the Collatz map attain almost bounded values}.
Forum of Mathematics, Pi \textbf{10}, e12, 2022.
\url{https://doi.org/10.1017/fmp.2022.8}.
Numbered references follow arXiv v7:
\url{https://arxiv.org/abs/1909.03562v7}.
\bibitem{Inselmann2024} M. Inselmann.
\emph{An Approximation of the Collatz Map and a Lower Bound for its
Average Total Stopping Time}.
Preprint, arXiv:2402.03276v3, 2024.
\url{https://arxiv.org/abs/2402.03276v3}.
\bibitem{Shaik2026} I. A. Shaik.
\emph{Polylogarithmic Descent for Almost All Collatz Orbits in Natural Density}.
Preprint, version 4.0.0, 6 September 2026.
\url{https://doi.org/10.5281/zenodo.22525724}.
Software: \url{https://doi.org/10.5281/zenodo.22525757}.
\bibitem{NashidaBarrier2026} H. Nashida.
\emph{A power-saving bound, uniform in the endpoint, for Collatz orbits
that stay above a fixed barrier}.
Preprint, revision r8, 27 September 2026.
\url{https://doi.org/10.5281/zenodo.22986671}.
\bibitem{NashidaGeneral2026} H. Nashida.
\emph{A power saving in natural density for generalized Collatz maps}.
Preprint, revision r5, 27 September 2026.
\url{https://doi.org/10.5281/zenodo.22986678}.
\bibitem{NashidaTerras2026} H. Nashida.
\emph{A strict improvement of the Terras-type bound for integers with
long Collatz stopping times}.
Preprint, revision r6, 27 September 2026.
\url{https://doi.org/10.5281/zenodo.22986642}.
\bibitem{MazurTao2026} L. Mazur and ProofAtlas.
\emph{Tao's Almost-Bounded Collatz Orbits}, Lean formalization, 2026.
\url{https://www.proofatlas.ai/formalizations/tao-almost-bounded-orbits/}.
Checked-source commit \texttt{d53c8de00056fb05999e44589f2399d7efa026e4}.
\bibitem{MazurNatural2026} L. Mazur.
\emph{Natural-Density Almost-Bounded Collatz Orbits in Logarithmic Time}.
Preprint, version 2, 21 July 2026.
\url{https://www.proofatlas.ai/formalizations/natural-density-log-time-collatz/}.
Checked-source package \texttt{ca3dd0d63920411213403092aecc6946619eb082}.
\bibitem{MazurPredecessors2026} L. Mazur.
\emph{Positive Lower Density of Collatz Predecessors}.
Preprint, version 1.0, 6 September 2026.
\url{https://www.proofatlas.ai/formalizations/positive-lower-density-collatz-predecessors/}.
Checked-source package \texttt{c290ac229b379d2b51026d3e86140496870d6f39}.
\end{thebibliography}
\end{document}
